\documentclass[11pt,a4paper]{article}
\usepackage[margin=25mm]{geometry}
\usepackage[T1]{fontenc}
\usepackage{hyperref}
\begin{document}
\title{B1 document hierarchy -- Phase 4 CJPE methodology}
\author{CS781 project documentation}
\date{Reconstructed from preserved source and run records}
\maketitle
\section{Problem statement and data}
Phase 4 Court Judgment Prediction predicts binary outcome 0/1 from legal case text. The preserved official inputs are Train/train.json, Train/dev.json and Test/test.json. Test labels are private. Original run records report 994 official dev labels and 1,517 test IDs; macro-F1 is the central metric.
\section{Architecture and mechanics}
B1 6-encoder/6-decoder original Phase 1 backbone; 512 tokens per chunk, 16 ordered chunks per case, overlap 100, tail sampling. Two document Transformer blocks followed by bidirectional GRU. Token representation uses a learned mixture of mean, attention and last content token; position enters before document aggregation. The final learned attention pools document states for two-class logits. Data flow: Legal case -> B1 6+6 backbone -> 512×16 windows -> Token evidence mix -> Two document Transformer blocks followed by bidirectional GRU -> Attention pool -> Binary ZIP.
\section{Scientific rationale}
B1: six-layer encoder and six-layer decoder from own Phase 1 weights; decoder representations supply token evidence. This experiment isolates the effect of its recorded adaptation or document-head mechanism while retaining the original task and export contract. Comparisons across different backbones, training budgets or data coverage remain descriptive rather than controlled causal estimates.
\section{Exact configuration, training progress and compute}
Head mode transformer\_gru; 6 encoder blocks, no FFN expansion in this scored configuration; window policy tail, 512 tokens per window, 16 windows, overlap 100. Backbone learning rate 1e-05, head learning rate 0.0003, accumulation 8, warmup updates 100, target maximum epochs 4. Latest recorded segment 4 in this package; selected step 6500, optimizer updates 8000, completed full epochs 2, seen training cases 63992. Selected checkpoint SHA-256 0626c85191eb9e74a2133152e336b2b0cf9ad3e9a082b315eb2d129f3ac10a90; peak GPU bytes 2520698368. The preserved lineage JSON traces the resumed optimizer/checkpoint state across segments; the ZIP is produced from a selected checkpoint, not an invented continuation.
\section{Exact prediction path}
After token evidence pooling, a width-256 projection plus normalized position features enters two document Transformer blocks (four heads, width-768 gated feed-forward inner layer), then a 128-unit-per-direction bidirectional GRU. A learned document attention pools states into two logits.
\section{Evaluation mechanics}
Official labeled dev cases 994; balanced class supports [497, 497]; selected-threshold confusion matrix [[375, 122], [322, 175]]; calibrated threshold 0.486789 and macro-F1 0.534473; default 0.5 macro-F1 0.499201. These are two evaluations of the same selected weights, with different binary decision thresholds. The 1,517-ID test prediction ZIP is a separate label-free artifact.
\section{Adaptation mechanism}
Staged fine-tuning starts at the task head and progressively unfreezes original encoder/decoder blocks. All optimizer parameter groups exist from the beginning; stage changes control which tensors receive gradients. Saved optimizer, scaler, RNG and position permit continuation between segments without reinitializing this head.
\section{Training design and checkpoint selection}
Staged B1 fine-tuning; document head mode `transformer\_gru`. Original B1 study notebook lineage `code/v1.ipynb` through `code/v4.ipynb`; latest source snapshot `code/source/b1\_study/`. Saved configs and terminal/run evidence are in `evidence/`.
\section{Training objective and validation boundary}
The B1 study uses cross-entropy for the binary label, AdamW with 0.01 weight decay, and staged unfreezing. Validation is performed on a fixed train-only holdout; checkpoint selection requires at least two epochs before patience can declare convergence. A time-budget stop instead records paused and a resumable checkpoint. The fixed 512-example train holdout selected model checkpoints at threshold 0.5. A separate, group-disjoint 512-example train calibration set selected the final decision threshold. The official 994-case dev evaluation used that threshold: macro-F1 0.534473; using 0.5 instead gave 0.499201. The split code excludes normalized-text duplicate groups across train, holdout, and calibration. The 1,517 test cases supplied no labels. A budget-paused run is not convergence.
\section{Measured results}
Version 4 run state **complete**; 8000 optimizer updates and 2 complete epochs, selected step 6500. Calibrated official 994-case dev macro-F1 **0.534473** at threshold 0.486789; default 0.5 dev macro-F1 0.499201. Selected checkpoint SHA-256 `0626c85191eb9e74a2133152e336b2b0cf9ad3e9a082b315eb2d129f3ac10a90`; verified 1,517-ID ZIP SHA-256 `eb70fc73838a3fa84a6bd6ef5eac03ce680f62ef93bf780e8a9775ea7a384d7a`. No finished Codabench score for this exact export is established by the supplied screenshot.
\section{Source and evidence trail}
B1 source and notebook versions are archived. The v7 source snapshot in code/source/ is the latest executable reference; for older v4 exports, the corresponding versioned notebook and run config identify the exact earlier run rather than implying that v7 code executed v4. The code SHA-256 manifest, original run records, selected-checkpoint identity, and verified prediction archive are indexed in evidence/README.md. This report is a later reconstruction of preserved evidence, not a claim that its companion notebook ran the original training.
\section{Prediction contract and verification}
The original prediction ZIP contains one root-level prediction.json with 1,517 unique official test IDs and integer labels restricted to 0 and 1. This package checks the ZIP CRC and SHA-256 against the archived run/verification record. Full model checkpoint tensors remain outside GitHub; their downloaded backup receipts and hashes live in the local Phase 4 study archives.
\section{Reproducibility and limits}
A budget-paused run is not converged; the target was four epochs with checkpoint selection and train-only calibration. The v7 source also defines gated depth expansion and an upper FFN option, but the recorded configs for these exports keep encoder\_depth=6 and expand\_ffn=false; those unlaunched options have no reported score. Original checkpoint bytes were downloaded and verified in the local archive, but are omitted from GitHub. Reports and companion notebook were created after training; original scores come from copied artifacts. Model weights and course data are preserved locally, not uploaded to GitHub.
\end{document}
